%!TEX root=./main.tex 

\subsection{Counting Execution Graphs: Hardness}

When counting execution graphs, we assume that all runs of $P$ are bounded (by some number $n$).
To simplify notation, we will assume in the following that we have ``compiled'' the number $n$ 
into the program (e.g., by unrolling loops an appropriate number of times) and 
therefore, every run of the program is already bounded by the size of the program. 
Thus, we will omit $n$ in our problem formulation and results, and assume $n = |P|$ 
bounds the number of events in any execution of $P$.
%
Our goal is to count the number $C(P)$ of maximal reachable execution graphs of $P$.
To this end, we define the following counting problem.

\begin{problem}[The Counting Problem]
    \label{prob1}
    Given as input a program $P$, output the number of maximal reachable execution graphs $C(P)$ in $\mathcal{T}(P)$.
\end{problem}

%To simplify notation, we will assume in the following that we have ``compiled'' the number
%$n$ into the program (e.g., by unrolling loops an appropriate number of times) and therefore, every execution of the program
%is already bounded by the size of the program.
%Thus, we will omit the input $n$ in our results and assume $n = |P|$.

We first show that the counting problem above is complete for the class $\sharpP$ (proof in \cref{theory:A}):
\begin{theorem}\label{thm:counting-problem-hardness}
   The Counting Problem is $\sharpP$-complete. 
   $\sharpP$-hardness holds even if either the number of threads or the number of shared locations is constant.
\end{theorem}

%\RM{check the sentence below:}
%In fact, the problem is $\sharpP$-hard already when either the number of threads or the number of locations (but not both) is fixed. \MZ{This is not correct. We showed that %if threads and locations are fixed but local register is not then the problem is also $\sharpP$-hard.
%We proved that when everything is constant it has a polynomial time algorithm because there is log-space transducer. }

Recall that $\sharpP$ is the complexity class associated with counting the number
of accepting computations of a non-deterministic Turing machine (see, e.g., \cite{aroraBarak}). 
In particular, the problem of counting the number of satisfying assignments of a SAT formula is 
$\sharpP$-complete. 
Hence, $\sharpP$-complete problems are at least as hard as NP-complete problems, which 
gives strong evidence against the existence of any polynomial-time
algorithms for Problem~\ref{prob1}. 
In light of this, we can relax the requirement of the Counting Problem 
so that we only need an approximate count, rather than an exact count.



\begin{problem}[The Approximate Counting Problem]
    \label{prob2}
    A randomized algorithm $A$ is an $(f(n), \rho)$-approximate counter if, for every program $P$ of size $n$, we have
    \[
    \operatorname{Pr}[C(P)/f(n) \leq A(P, \rho) \leq C(P)f(n)] \geq 1 - \rho
    \]
%    Given as input a program $P$ of size $n$, output a number $N$ such that $N/2^{n^{1-\epsilon}} \le C(P) \le N \times 2^{n^{1-\epsilon}}$
%    for some constant $\epsilon > 0$ independent of the input.
\end{problem}

Note that since $C(P)$ is always at most an exponential function in the size of $P$,
the least non-trivial approximation that we can ask for is an approximation
with a sub-exponential ratio. 
However, even for this very weak approximation ratio, we show that it is highly unlikely that a (randomized) polynomial-time algorithm exists, by means of the following theorem
(proof in \cref{theory:B}):

%\begin{theorem}
 %Let $d>0$ be any constant, and let $C$ be the function that computes Problem \ref{prob1}. There is no polynomial algorithm which always output an approximation $\mathcal{A}$ for $C(P)$ such that
  % $\mathcal{A}/2^{n^{d}}\leq  C(P)\leq \mathcal{A}*2^{n^{d}}$ for every concurrent program $P$ unless P=NP.
%\end{theorem}
\begin{theorem}\label{thm:approximate-counting-problem-hardness}
 Unless RP = NP, there is no randomized polynomial-time $(2^{n^{1-\epsilon}}, \rho)$-approximate 
 counter for $\epsilon > 0$ independent of the input.
\end{theorem}

Here RP is the class of problems which admit randomized polynomial-time algorithms;
(see \cite{aroraBarak}). Hence, this theorem puts hard limits on the Approximate Counting Problem. 
Furthermore, unless P = NP, a similar argument shows that there is no deterministic polynomial-time
algorithm that gives a sub-exponential approximation to $C(P)$. 
%(We defer the proofs of both of these theorems to ~\cref{theory}).

In light of these hardness results, we now present some polynomial-time randomized algorithms that offer an \emph{unbiased estimate} of $C(P)$, i.e., the \emph{expected value} returned by them will be $C(P)$. 



% In the rest of the paper, we begin by presenting two Monte Carlo estimation methods that are provably unbiased, offering principled solutions to the estimation problem. We then establish the hardness of exact counting to provide a clearer theoretical understanding of the problem’s complexity. Finally, we evaluate the practical effectiveness of our proposed estimators on a set of concurrent programs by comparing them against exact counts obtained via a model checker.


%\begin{problem}[Approximate Counting with Additive Errors]
%    \label{prob3}
%\RM{TODO FIX}
%    Given as input a program $P$ of input size $n$, output a number $N$ such that $| C(P) - N| \le poly(n, 1/\epsilon)$
%    for some constant $\epsilon > 0$ independent of the input.
%\end{problem}
